\section{Simulator, data, replay and contamination}\label{app:simdata}

\subsection{Simulator and cost model}

The data store returns rows at or before the decision index. The agent observes the close at bar $t$ and its order fills at that same close, with costs applied to the fill price. That close-to-close convention is idealized; the direction and size of a decision delay's effect on each entrant are not established, and a delay sensitivity is planned. Execution charges a per-trade fee, seeded slippage, a size-dependent rebalancing cost that scales with the square root of the trade's value as a fraction of the agent's own net asset value, and financing on gross exposure above one times net asset value. A per-step turnover cap can limit each trade to a fixed fraction of net asset value. Both the rebalancing cost and the cap are measured against the agent's book: the datasets carry no volume, so neither is a share of market volume, capacity and volume-based market impact are not modelled, and a small book and a large one pay the same rate. Four named cost profiles ship, and each is reported under the name used throughout this paper rather than under its Rust identifier: frictionless (\texttt{CostProfile::None}), typical (\texttt{CostProfile::Typical}), stressed (\texttt{CostProfile::WorstCase}) and realistic (\texttt{CostProfile::Realistic}). The typical profile is 2 basis points fee, 3 slippage, 50 rebalancing cost and 5 financing; stressed is 10, 15, 150 and 20 with a per-step turnover cap of 10 percent of net asset value and a two-bar delay, which the backtest driver does not apply; frictionless charges none of them. Realistic is typical plus seeded fill delay, partial fills and queue-position slippage. \Cref{sec:external} sweeps frictionless, typical and stressed; \cref{sec:seedleg} compares typical with realistic on three daily datasets; every other experiment uses the typical profile. The seeded slippage is what makes different execution seeds produce different runs and gives $\passk$ something to measure.

\subsection{Data}\label{sec:data}

Nine frozen datasets supplied the historical benchmark evidence, each with a SHA-256 sidecar and a deterministic fetch or derivation script. The source requests needed no API key, but that access fact did not grant redistribution permission; \cref{app:datasheet} records the unresolved rights blocker. \Cref{tab:data} lists the exact historical bytes and date ranges. Before any agent is scored on a dataset, the dataset is scored: a battery of stylized facts after \citet{cont2001empirical} tests excess kurtosis, volatility clustering, gain-loss asymmetry and aggregational Gaussianity, and time-reversal asymmetry after \citet{zumbach2009time}, and issues a verdict. Seven of the nine pass. Weekly US indices fail the aggregation test because 522 bars is too short for it, and FX majors fail the time-reversal asymmetry test, which does not detect the asymmetry at its threshold on these daily noon buying-rate series. Both ship with the verdict recorded. Apart from crypto spot, every series is a price level rather than a tradable total-return position: the equity series are price indices without dividends, the FX series are spot rates without carry and the oil series are spot prices without roll yield (\cref{sec:stats}).

Two datasets carry caveats the scripts document. The commodities file keeps the WTI close of $-36.98$ on 20 April 2020 as published, which dominates its moments, so every commodities statistic in the paper inherits that outlier; a start date of May 2020 is the documented opt-out and removes it. The rates file's close column is a yield in percent, not a price. The simulator nevertheless treats that column as a price level, so its percentage changes are neither Treasury total returns nor the P\&L of a duration-specified bond position. That row is a numerical stress case for the protocol on a one-symbol universe and must not support an economic performance claim: cross-sectional momentum is undefined there, and the luck floor requires the single-symbol randomization of \cref{sec:falsify} to be a control at all. \Cref{sec:limitations} states the coverage gaps.

\subsection{Replay and recompute}\label{sec:replay}

A captured trajectory contains the agent's raw decisions and a schema-v2 \texttt{TrajectoryContract}, never a return or self-reported statistic. The contract names the dataset and cost-model digests, engine version, exact ordered windows and seeds and, for command-line capture, the runner executable digest. Strict verification refuses an empty artifact, a missing or old contract, any identity mismatch, duplicate or out-of-range geometry, a run matrix other than the complete window-major Cartesian product, an incorrect number or order of decision steps, and an observation date different from the frozen bar at that step. A separate verifier then replays the decisions through the same engine against the frozen dataset and recomputes the score. Capture and replay share the engine path, and a test asserts that doubling every order in a captured trajectory recomputes to different returns. The default command is strict; \texttt{--allow-unbound-trajectory} is an explicit diagnostic escape hatch for legacy or cross-version regrading and cannot be mistaken for the normal verifier. An artifact therefore cannot substitute self-reported returns or silently replay against different recorded conditions. What it cannot do is prove that the committed strategy produced those decisions; the trajectory contract carries no entrant commitment, so that linkage remains separate from the commitment in \cref{sec:attest}. The operator's rescore and regrade commands compose with this verifier rather than replacing it: a rescore recomputes from a declared bundle whose manifest is the whole read set, and a regrade emits a receipt that links a frozen trajectory, hashed from its bytes, to a superseding evaluator without publishing a figure, and without replacing a source listed as frozen published evidence.


\subsection{Contamination}\label{sec:contamination}

Two defenses ship. A held-out dataset can be sealed under a published commitment to its hash, so the benchmark can show later that the data a model was scored on is byte-identical to the data it committed to before scoring. As in \cref{sec:attest}, the commitment binds bytes and not chronology. A canary token with the standard training-corpus warning is embedded beside it, so post-hoc detection of a verbatim copy of the scenarios in a model's output is possible. The signal is positive-only and literal-bytes-only: a hit is evidence of contamination, and a miss is not evidence of its absence, because a paraphrased, compressed, encoded or model-internalized copy carries none of the protected bytes. Separately, a masked view of any dataset renames symbols and dates to opaque identifiers, which defeats an agent that pattern-matches memorized tickers.

\subsection{The units correction, in detail}\label{app:unitsfix}

The fix, shipped in v0.3.0, leaves the per-period kernel untouched and gives the configuration a periods-per-year field so that $\sigma_{\text{trials}}$ is stated annualized and converted once at the point of use. A test asserts the conversion is applied exactly once on the configured path and never on the measured path. Another asserts that an SPX-like series is eligible under an annualized prior of $0.1$ at 252 periods per year and ineligible under the same prior at one period per year, which reproduces the old units exactly. A test that the same series clears the original $0.5$ prior was not written, because it cannot: $0.5$ annualized at 50 trials is a benchmark of $1.14$ annualized, and the index sits below it at any track length. That is the correct verdict for that prior, and the benchmark now states it in units a reader can check.
